%\documentclass[aspectratio=169, handout]{beamer}    % For handouts.
\documentclass[aspectratio=169]{beamer}             % For slides.

\input{../../../../hh_bpa_forel-preamble}

\title{FMTS mätteknik HT2026 -- Dag 1}
\author{Pererik Andreasson \\ Struan Gray }
\date{\today{}}

\usepackage{circuitikz}
\usepackage{pgfplots}
\pgfplotsset{compat=1.18}

\begin{document}

%%% Title slide: different footer, large logo
{
% no page #, no logo on title page
\setbeamertemplate{footline}{}
\setbeamertemplate{logo}{\titlelogos}
\begin{frame}
\vfill
  \begin{block}
    \maketitle
  \end{block}
\end{frame}
}

% Preamble:

\begin{frame}{\(\sin x\)}
    \centering
    \begin{tikzpicture}[
        line width=1.2pt,
        every node/.style={font=\Large}
    ]
        \coordinate (A) at (0,0);
        \coordinate (B) at (6,0);
        \coordinate (C) at (6,3.5);

        \draw (A) -- (B)
                  -- node[right] {$a$} (C)
                  -- node[above left] {$c$} (A);

        % Right angle and angle x
        \pic[draw, angle radius=0.35cm]
            {right angle=A--B--C};
        \pic[draw, "$x$", angle radius=1.1cm, angle eccentricity=1.3]
            {angle=B--A--C};
    \end{tikzpicture}

    \bigskip
    {\LARGE \(\displaystyle \sin x=\frac{a}{c}\)}
\end{frame}


\begin{frame}{The unit circle}
    \centering
    \begin{tikzpicture}[
        scale=2.3,
        line width=1.2pt,
        every node/.style={font=\large}
    ]
        % Same triangle proportions as before
        \pgfmathsetmacro{\ang}{atan(3.5/6)}

        \coordinate (O) at (0,0);
        \coordinate (P) at (\ang:1);
        \coordinate (B) at ({cos(\ang)},0);

        % Coordinate axes and unit circle
        \draw[->, thin, gray] (-1.25,0) -- (1.3,0);
        \draw[->, thin, gray] (0,-1.2) -- (0,1.25);
        \draw (O) circle[radius=1];

        \node[below left] at (O) {$0$};
        \node[below right] at (1,0) {$1$};
        \node[above left] at (0,1) {$1$};

        % Inscribed right triangle
        \draw (O) -- (B)
                  -- node[right] {$\sin x$} (P)
                  -- node[above left] {$1$} (O);

        \pic[draw, angle radius=0.18cm]
            {right angle=O--B--P};
        \pic[draw, "$x$", angle radius=0.7cm,
             angle eccentricity=1.35]
            {angle=B--O--P};

        \fill (P) circle[radius=0.025];
    \end{tikzpicture}

    \medskip
    {\Large \(\displaystyle \sin x=\frac{a}{1}=a\)}
\end{frame}


\begin{frame}{Seriekopplade komponenter.}
\begin{center}
\resizebox{0.4\textwidth}{!}{
\begin{tikzpicture}[american]
	\draw (0, 0) to[sV, v<=\qty{2}{V}] (4, 0)
		to[short] (4, 2)
		to[C, v=$U_C$] (2, 2)
		to[R, v=$U_R$] (0, 2)
		to[short] (0, 0) ;
\end{tikzpicture}
} % end resizebox
\end{center}
\end{frame}

\end{document}